% ============================================================
% Round 01: SFT
% ============================================================

\subsection{Round 01：SFT}

\subsubsection{本轮问题}

\begin{conceptbox}
\textbf{学习目标：} 先理解 LLM 后训练链条里的第一步：SFT。\\
\textbf{主线位置：}
\[
\text{SFT} \to \text{Reward Model} \to \text{PPO / RLHF} \to \text{DPO} \to \text{GRPO / RLVR}
\]
\end{conceptbox}

\subsubsection{后训练在解决什么}

预训练模型学到的是：

\[
\text{给定上下文，预测下一个 token}
\]

它会语言、知识和模式，但不一定会按人类想要的方式工作。例如：

\begin{itemize}[leftmargin=1.5em]
  \item 不一定听指令；
  \item 不一定给结构化答案；
  \item 不一定安全；
  \item 不一定偏向高质量推理；
  \item 不一定知道什么回答更好。
\end{itemize}

所以后训练的目标是：

\[
\text{Base Model} \to \text{Assistant / Reasoner / Agent}
\]

\begin{summarybox}
\textbf{最小定义：} Post-training = 用额外信号重塑模型的输出分布。
\end{summarybox}

\subsubsection{SFT 是什么}

SFT = Supervised Fine-Tuning，即监督微调。

数据形式：

\begin{center}
{\small
\begin{tabular}{p{3cm}p{10cm}}
\hline
\textbf{输入 $x$} & 请解释什么是 GRPO \\
\textbf{目标答案 $y^\ast$} & GRPO 是一种不使用 critic 的 policy optimization 方法…… \\
\hline
\end{tabular}
}
\end{center}

训练目标：

\[
\mathcal{L}_\text{SFT} = - \log \pi_\theta(y^\ast \mid x)
\]

展开到 token 级别：

\[
\mathcal{L}_\text{SFT}
=
-\sum_{t=1}^{T}
\log \pi_\theta(y_t^\ast \mid x, y_{<t}^\ast)
\]

含义是：让模型在看到 prompt $x$ 时，更倾向于生成人工写好的目标答案 $y^\ast$。

\subsubsection{SFT 的直觉}

Base model 的行为更像：

\[
\text{看到一句话，继续往下写}
\]

SFT 之后，模型更像：

\[
\text{看到一个任务，给出符合格式和意图的回答}
\]

因此 SFT 主要学习：

\begin{itemize}[leftmargin=1.5em]
  \item 指令格式；
  \item 回答风格；
  \item 任务模板；
  \item 基础对齐行为。
\end{itemize}

\begin{intubox}
\textbf{例子：} 对于“总结这篇论文”这个 prompt，SFT 前模型可能只是续写论文风格文本；SFT 后模型更可能输出“背景、方法、实验、结论”这样的任务型结构。
\end{intubox}

\subsubsection{SFT 的本质}

SFT 本质上是：

\[
\text{Behavior Cloning / Imitation Learning}
\]

它不是在判断哪个答案更好，而是在模仿给定的标准答案。因此能力上限高度依赖数据质量：

\begin{center}
{\small
\begin{tabular}{p{4cm}p{8.5cm}}
\hline
\textbf{数据质量} & \textbf{结果} \\
\hline
高质量 SFT 数据 & 行为更稳定，格式更清晰，任务遵循更好 \\
低质量 SFT 数据 & 学会啰嗦、幻觉、格式污染或错误偏好 \\
\hline
\end{tabular}
}
\end{center}

\subsubsection{SFT 的局限}

\begin{enumerate}[leftmargin=1.5em]
  \item \textbf{不能区分多个答案的相对好坏。} 如果两个答案都正确，但一个更清晰，SFT 只会模仿数据里的那个，不天然知道哪个更好。
  \item \textbf{不会主动探索更优答案。} SFT 只学已有示范，不会自己采样、比较和改进。
  \item \textbf{强依赖人工答案质量。} 示范答案差，模型会稳定地学到差行为。
  \item \textbf{对推理能力提升有限。} SFT 可以教模型输出推理格式，但不一定真正提升找对答案的能力。
\end{enumerate}

这就是后面需要 Reward Model、PPO、DPO、GRPO 的原因：它们试图从“模仿标准答案”进一步走向“优化更好的答案”。

\subsubsection{本轮最小结论}

\begin{summarybox}
\textbf{SFT = 给模型标准答案，让它模仿。}\\
核心公式：
\[
\mathcal{L}_\text{SFT} = - \log \pi_\theta(y^\ast \mid x)
\]
核心局限：SFT 只能学数据里怎么答，不能直接学哪个答案更好。
\end{summarybox}

\subsubsection{本轮检查题}

\begin{enumerate}[leftmargin=1.5em]
  \item SFT 为什么能让 base model 变得更像 assistant？
  \item SFT 的 loss 在优化什么？
  \item 为什么 SFT 之后还需要 RLHF / DPO / GRPO？
\end{enumerate}
